\documentclass{article} 

\usepackage{mdframed}
\usepackage{listings}
\usepackage{hyperref}
\usepackage{amsfonts}
\usepackage{amsmath}

\title{CSE160 Homework 2}
\author{Elijah Hantman}

%\renewcommand{\spn}[1]{\langle #1 \rangle}


\begin{document}
    \maketitle

    {\large Part 1}
    \begin{enumerate}
        \item Using the labels Translate, Rotation, Scale, or Reflect, label the transformation of
            point p when multiplied by matrix $M(p' = Mp)$:

            \begin{enumerate}
                \item 
                    \begin{displaymath}
                        M = 
                        \begin{bmatrix}
                            0 & -1\\ 
                            1 & 0
                        \end{bmatrix}
                    \end{displaymath}

                    \begin{mdframed}
                        This is a rotation by 90 degrees.
                    \end{mdframed}
                \item 
                    \begin{displaymath}
                        M =
                        \begin{bmatrix}
                            \frac{\sqrt{2}}{2} & \frac{\sqrt{2}}{2}\\ 
                            -\frac{\sqrt{2}}{2} & \frac{\sqrt{2}}{2}
                        \end{bmatrix}
                    \end{displaymath}
                    \begin{mdframed}
                       This is another rotation, by 45 degrees. 
                    \end{mdframed}
                \item 
                    \begin{displaymath}
                        M =
                        \begin{bmatrix}
                            2 & 0 & 0 & 0\\
                            0 & 4 & 0 & 0\\
                            0 & 0 & 2 & 0\\
                            0 & 0 & 0 & 1\\
                        \end{bmatrix}
                    \end{displaymath}
                    \begin{mdframed}
                        This is a non uniform scaling.
                    \end{mdframed}
                \item 
                    \begin{displaymath}
                        M = 
                        \begin{bmatrix}
                            -1 & 0\\ 
                            0 & 1
                        \end{bmatrix}
                    \end{displaymath}
                    \begin{mdframed}
                        This is a reflection about $e_2$,
                        meaning a vector $v = ae_1 + be_2$
                        will be mapped to $v = -ae_1 + be_2$.
                    \end{mdframed}

                \item 
                    \begin{displaymath}
                        M =
                        \begin{bmatrix}
                            1 & 0 & 3\\ 
                            0 & 1 & 2\\
                            0 & 0 & 1
                        \end{bmatrix}
                    \end{displaymath}
                    \begin{mdframed}
                        In a 2D projective space, this would
                        correspond to a translation by
                        $(3,2)$.
                    \end{mdframed}
            \end{enumerate}

        \item Find the matrix for a $120\circ$ rotation about the axis defined by the vector $r = (1,1,0)$.
            \begin{mdframed}
                The matrix can be derived via the product of a rotation to put the axis onto one of
                the bases, then a basis aligned rotation, followed by the inverse of the original rotation.

                Using a tool to calculate the rotation matrix we get:

                \begin{displaymath}
                    M = 
                    \begin{bmatrix}
                        0.2500000&  0.7500000&  0.6123725\\
                        0.7500000&  0.2500000& -0.6123725\\
                        -0.6123725&  0.6123725& -0.5000000 
                    \end{bmatrix}
                \end{displaymath}

                Alternative methods involve using quaternions,
                then converting to a matrix, or you could decompose
                the rotation using euler angles. Ie: first doing
                a rotation to put the axis on the basis, then doing
                a rotation, then undoing it.
            \end{mdframed}

        \item Match the following 2D homogeneous matrices to the transformations in the image:
            \begin{mdframed}
               1 goes with matrix b.\\ 
               2 goes with matrix a\\
               3 goes with matrix c
            \end{mdframed}

        \item Describe a sequence of Translate(x,y), Rotate(degrees), and/or Scale (by a value
            of ) that when multiplied describe the below transformation to go from the ’Before’
            to ’After’.

            \begin{mdframed}
                Translate (-1, 0)\\
                Rotate(45)
                
                \vspace{10pt}

                I can't tell if the after image has
                a scale applied, so I didn't include that,
                but if I did, it would have to happen first.
            \end{mdframed}

            \break
        \item 
            \begin{enumerate}
                \item Describe (using a sequence of (3 x 3) matrix multiplications) the transformations
needed to transform the triangle from A to B in the figure below:

                \begin{mdframed}
                    \begin{displaymath}
                        \begin{bmatrix}
                           x'\\ 
                           y'\\
                           z'
                        \end{bmatrix}
                        =
                        \begin{bmatrix}
                            1 & 0 & -7\\ 
                            0 & 1 & 5\\
                            0 & 0 & 1
                        \end{bmatrix}
                        \begin{bmatrix}
                            2 & 0 & 0\\ 
                            0 & -1 & 0\\
                            0 & 0 & 1
                        \end{bmatrix}
                        \begin{bmatrix}
                         x\\   
                         y\\
                         z
                        \end{bmatrix}
                    \end{displaymath}

                    Something to note is that I do this
                    in 2 steps. First I scale and reflect
                    the triangle, then I translate it into
                    position. I also do not forget that scaling
                    far from the origin causes the position to
                    shift.
                \end{mdframed}

            \item Using the transformations found in part (a), multiply the following points with
                the matrix.

                \begin{displaymath}
                    p_1 = \begin{bmatrix}
                       4\\ 
                       1
                    \end{bmatrix}
                    \quad
                    p_2 =
                    \begin{bmatrix}
                       4\\ 
                       3
                    \end{bmatrix}
                \end{displaymath}

                \begin{mdframed}
                    \begin{displaymath}
                        p_1 = (4,1)
                    \end{displaymath} 
                    \begin{displaymath}
                        \to (8,-1)
                    \end{displaymath}
                    \begin{displaymath}
                        \to (1, 4)
                    \end{displaymath}

                    \begin{displaymath}
                        p_2 = (4,3)
                    \end{displaymath}
                    \begin{displaymath}
                        \to (8, -3)
                    \end{displaymath}
                    \begin{displaymath}
                        \to (1, 2)
                    \end{displaymath}
                \end{mdframed}
            \end{enumerate}
    \end{enumerate}

    \break
    {\large Part 2}
    \begin{enumerate}
        \item Define uniform, attribute and varying variables
            \begin{mdframed}
                Uniform variables are variables which are
                constant between vertex shader invocations.

                On the other hand attribute and varying variables
                contain the data from vertices, and interpolated
                vertex shader outputs respectively.

                In more modern GLSL uniforms still exist since
                the idea of simply making a single value accessible
                across all shader executions is very useful.

                On the other hand in order to make the API more
                descriptive attribute and varying variables
                are simply labeled as "in" and "out" variables.

                This means that the variables are more clearly
                defined in regards to the shader they appear in.
                For example, a vertex shader will output to
                an "out" variable, which would be labeled
                varying in GLSL 1.

                And a fragment shader will take as input an
                "in" variable rather than a varying variable.
            \end{mdframed}

        \item What type is the shader and describe its function?
            \begin{enumerate}
                \item

                    \begin{lstlisting}[language=C]
attribute vec2 a_Position;
main() {
  gl_Position = vec4(0.5 * a_Position, 0.0, 1.0);
} 
                    \end{lstlisting}

                    \begin{mdframed}
                        This is a vertex shader. We can tell by
                        the use of "attribute" to declare vertex
                        data, as well as "gl\_Position" which
                        is a vertex shader builtin which
                        is used to set the position of the
                        vertex for rasterization.
                    \end{mdframed}

                \item 
                    \begin{lstlisting}[language=c]
main() {
  gl_FragColor = vec4(0.0, 0.0, 1.0, 1.0);
} 
                    \end{lstlisting}

                    \begin{mdframed}
                       This is a fragment shader. This
                       is obvious by the usage of
                       "gl\_FragColor" as well as the
                       lack of attribute declarations.
                    \end{mdframed}
            \end{enumerate}

            \break
        \item Using linear interpolation, calculate the color of the following points that are on the
            line defined by the 2D points p1 = [3,5] with color [255,0,0] and p2 = [10, 5] with
            color [0,0,255].
            \begin{enumerate}
                \item $[4.8, 5]$
                    \begin{mdframed}
                        \begin{displaymath}
                            [4.8, 5] = t[10.5] + (1-t)[3, 5]
                        \end{displaymath}
                        \begin{displaymath}
                            [4.8, 5] = tp_2 + p_1 - tp_1
                        \end{displaymath}
                        \begin{displaymath}
                            [4.8, 5] = t(p_2 - p_1) + p_1
                        \end{displaymath}
                        \begin{displaymath}
                            [1.8, 0] = t(p_2 - p_1)
                        \end{displaymath}
                        \begin{displaymath}
                            [1.8, 0] = t([7,0])
                        \end{displaymath}
                        \begin{displaymath}
                            t \approx 0.25
                        \end{displaymath}

                        Now we lerp the colors.

                        \begin{displaymath}
                            c = t(c_2 - c_1) + c_1
                        \end{displaymath}
                        \begin{displaymath}
                            c = t([-255, 0, 255]) + [255, 0, 0]
                        \end{displaymath}
                        \begin{displaymath}
                            c \approx [192, 0, 64]
                        \end{displaymath}
                    \end{mdframed}

                \item $[7.5, 5]$
                    \begin{mdframed}

                        Since most of the steps in the previous
                        problem didn't require using the
                        point we can skip most of the steps.

                        \begin{displaymath}
                            [4.5,0] = t([7,0])
                        \end{displaymath}
                        \begin{displaymath}
                            t \approx 0.64
                        \end{displaymath}

                        Now to lerp the colors
                        \begin{displaymath}
                            c = t([-255, 0, 255]) + [255, 0, 0]
                        \end{displaymath}
                        \begin{displaymath}
                            c \approx [92,0,163]
                        \end{displaymath}
                    \end{mdframed}
                    \break
                \item $[9, 5]$
                    \begin{mdframed}
                        \begin{displaymath}
                            [6, 0] = t([7,0])
                        \end{displaymath}
                        \begin{displaymath}
                            t \approx 0.86
                        \end{displaymath}

                        Next the colors:

                        \begin{displaymath}
                            c = t([-255, 0, 255]) + [255, 0, 0]
                        \end{displaymath}
                        \begin{displaymath}
                            c \approx [36,0,219]
                        \end{displaymath}
                        
                    \end{mdframed}
            \end{enumerate}
        \item Using barycentric coordinates, calculate the color for each of the following points
            inside the triangle defined by points $p_1 = [0,0,0]$, $p_2 = [3,5,0]$ and $p_3 = [6,0,0]$ with
            respective RGB colors $c_1 = [255, 0, 0]$, $c_2 = [0, 255, 0]$, $c_3 = [0, 0, 255]$.

            \begin{enumerate}
                \item $[1,1,0]$
                    \begin{mdframed}
                        To calculate barycentric coordinates you do,
                        \begin{displaymath}
                            v_1 = p_2 - p_1
                        \end{displaymath}
                        \begin{displaymath}
                            v_2 = p_3 - p_1
                        \end{displaymath}
                        \begin{displaymath}
                            v_3 = [1,1,0] - p_1
                        \end{displaymath}
                        \begin{displaymath}
                            A = |v_1 \times v_2|
                        \end{displaymath}
                        \begin{displaymath}
                            b_2 = |v_3 \times v_2|
                        \end{displaymath}
                        \begin{displaymath}
                            b_3 = |v_3 \times v_1|
                        \end{displaymath}
                        \begin{displaymath}
                            b_1 = 1 - (b_2 + b_3)
                        \end{displaymath}

                        Then to get the color we do
                        $b_1c_1 + b_2 c_2 + b_3 c_3$.

                        \begin{displaymath}
                            b = [11/15, 1/5, 1/15]
                        \end{displaymath}
                        \begin{displaymath}
                            c = [187, 51, 17]
                        \end{displaymath}
                    \end{mdframed}

                \item $[3,4,0]$
                    \begin{mdframed}
                        \begin{displaymath}
                            b = [1/10, 4/5, 1/10]
                        \end{displaymath}
                        \begin{displaymath}
                            c = [26, 204, 26]
                        \end{displaymath}
                    \end{mdframed}
                \item $[5, 0.25, 0]$.
                    \begin{mdframed}
                        \begin{displaymath}
                            b = [0.14, 0.05, 0.81]
                        \end{displaymath}
                        \begin{displaymath}
                            c = [36, 13, 207]
                        \end{displaymath}
                    \end{mdframed}
            \end{enumerate}
        \item Calculate the RGB color value for point P in the diagram below using bilinear
            interpolation. Show the intermediate values that you compute for interpolating
            along one of the two axes. The colors of the points A, B, C and D are the following:

            \begin{displaymath}
                A = RGB(0,0,0) \qquad B = RGB(200, 200, 200)
            \end{displaymath}
            \begin{displaymath}
                C = RGB(200, 0, 0) \qquad D = RGB(200, 0, 0)
            \end{displaymath}
            \begin{displaymath}
                A = (14,21) \qquad B = (15,21)
            \end{displaymath}
            \begin{displaymath}
                C = (14,20) \qquad D = (15,20)
            \end{displaymath}

            \begin{displaymath}
                P = (14,75, 20.5)
            \end{displaymath}

            \begin{mdframed}
                So we can instantly see that the lerp
                $C \to D$ will always give the result
                $RGB(200, 0, 0)$. This follows since
                $C = D$.

                Next we need to interpolate $A \to B$.

                \begin{displaymath}
                    14.75 = t(15 - 14) + 14
                \end{displaymath}
                \begin{displaymath}
                    0.75 = t
                \end{displaymath}

                Then the color.

                \begin{displaymath}
                    c = 0.75(RGB(200,200,200) - RGB(0,0,0)) + 
                    RGB(0,0,0)
                \end{displaymath}
                \begin{displaymath}
                    c = RGB(150, 150, 150)
                \end{displaymath}

                Next we interpolate vertically. Since the quad
                is axis aligned, we can just use the $y$ coordinate.

                \begin{displaymath}
                    20.5 = t(21 - 20) + 20
                \end{displaymath}
                \begin{displaymath}
                    0.5 = t
                \end{displaymath}

                Next the interpolation

                \begin{displaymath}
                    c = 0.5(RGB(150, 150, 150) - RGB(200, 0, 0)) + RGB(200, 0, 0)
                \end{displaymath}
                \begin{displaymath}
                    c = 0.5(RGB(-50, 150, 150)) + RGB(200, 0, 0)
                \end{displaymath}
                \begin{displaymath}
                    c = RGB(175, 75, 75)
                \end{displaymath}
            \end{mdframed}

        \item Given a polygon with vertices P1, P2, P3, each represented by a 3D vector $\begin{bmatrix} x & y & z\end{bmatrix}$,
            what formula calculates the face’s normal?

            \begin{mdframed}
                The normal is the vector facing directly perpendicular to the face.

                Assuming the winding order is $P1, P2, P3$.

                \begin{displaymath}
                    v_1 = P2 - P1
                \end{displaymath}
                \begin{displaymath}
                    v_2 = P3 - P1
                \end{displaymath}
                \begin{displaymath}
                    p = v_1 \times v_2
                \end{displaymath}
                \begin{displaymath}
                    n = \frac{p}{|p|}
                \end{displaymath}

                Since the cross product is defined as the vector perpendicular to
                $v_1$ and $v_2$. However the magnitude is the product 
                $|v_1| |v_2| sin(\theta)$. So to get a normal with unit length
                we need to normalize the vector.
            \end{mdframed}
    \end{enumerate}

    {\large Part 3}
    \begin{enumerate}
        \item Given a point $p = (x, y, z)$ and the point $p' = (x', y', z')$ where $p' = (x', y', z')$ is the point $p = (x, y, z)$ trans-
            lated by $T_x, T_y, T_z$:

            \begin{enumerate}
                \item Show the equations to compute $p'$ from $p$
                    \begin{mdframed}
                        \begin{displaymath}
                            x' = x + T_x
                        \end{displaymath}
                        \begin{displaymath}
                            y' = y + T_y
                        \end{displaymath}
                        \begin{displaymath}
                            z' = z + T_z
                        \end{displaymath}
                    \end{mdframed}
                \item Show the $4\times 4$ transformation matrix which takes
                    $p$ to $p'$.
                    \begin{mdframed}
                        \begin{displaymath}
                            M =
                            \begin{bmatrix}
                                1 & 0 & 0 & T_x\\ 
                                0 & 1 & 0 & T_y\\
                                0 & 0 & 1 & T_z\\
                                0 & 0 & 0 & 1 
                            \end{bmatrix}
                        \end{displaymath}
                    \end{mdframed}
            \end{enumerate}
            \break
        \item Given a point $p(x, y, z)$ and the point $p'(x', y', z')$, where $p'(x', y', z')$ is the $\beta$ degree rotated point
            of $p(x, y, z)$ around the z-axis:

            \begin{enumerate}
                \item Show the equation to compute $p'$ from $p$.
                    \begin{mdframed}
                       \begin{displaymath}
                           x' = x cos(\beta) - y sin(\beta)
                       \end{displaymath} 
                       \begin{displaymath}
                           y' = x sin(\beta) + y cos(\beta)
                       \end{displaymath}
                       \begin{displaymath}
                           z' = z
                       \end{displaymath}
                    \end{mdframed}
                \item Show the $4 \times 4$ transformation matrix.
                    \begin{mdframed}
                        \begin{displaymath}
                            M =
                            \begin{bmatrix}
                                cos(\beta) & -sin(\beta) & 0 & 0\\
                                sin(\beta) & cos(\beta) & 0 & 0\\
                                0 & 0 & 1 & 0\\
                                0 & 0 & 0 & 1
                            \end{bmatrix}
                        \end{displaymath}
                    \end{mdframed}
            \end{enumerate}
        \item Given a point $p(x, y, z)$ and the point $p'(x', y', z')$, where $p'(x', y', z')$ is the point $p(x, y, z)$ scaled
            by $S_x, S_y, S_z$:

            \begin{enumerate}
                \item Show the equation to compute $p'$ from $p$.
                    \begin{mdframed}
                        \begin{displaymath}
                            x' = x S_x
                        \end{displaymath}
                        \begin{displaymath}
                            y' = y S_y
                        \end{displaymath}
                        \begin{displaymath}
                            z' = z S_z
                        \end{displaymath}
                    \end{mdframed}
                \item Show the $4 \times 4$ transformation matrix.
                    \begin{mdframed}
                        \begin{displaymath}
                            M = 
                            \begin{bmatrix}
                                S_x & 0 & 0 & 0\\
                                0 & S_y & 0 & 0\\
                                0 & 0 & S_z & 0\\
                                0 & 0 & 0 & 1
                            \end{bmatrix}
                        \end{displaymath}
                    \end{mdframed}
            \end{enumerate}

        \item What is a model matrix? Write a simple vertex shader that uses a model matrix. Explain your
            shader.

            \begin{mdframed}
                A model matrix is a matrix which transforms objects from
                a local space where they contain their own origin, to world
                space where all objects are positioned relative to each other.

                In practice this means a model matrix is unique per object,
                and can do transformations before the camera and view matrices
                transform the scene into clip space.


                \begin{lstlisting}[language=C]
#version 300 es
in vec3 pos;             

uniform mat4 model;

main() {
    gl_Position = model * vec4(pos, 1.0);
}
                \end{lstlisting}

                This is a super simple vertex shader. It uses
                GLSL 2 which is why it has the version declaration
                at the top.

                Basically it reads in a vec3 from the vertex
                data, and a uniform model matrix. Then
                it sets the position to the model matrix
                applied to the position via left multiplication.

                This shader is one what would be used via
                a draw call per object with the model matrix
                being set per object. I also use a common
                pattern of converting the position to a vec4,
                rather than storing it as a vec4 since we
                always want the w component to be 1.
            \end{mdframed}
    \end{enumerate}

\end{document}
